\documentclass[a4paper,12pt]{article}

\usepackage[serbian]{babel}
\usepackage[T1]{fontenc}
\usepackage[utf8]{inputenc}

\usepackage{lmodern}
\usepackage{microtype}
\usepackage[a4paper,margin=2.5cm]{geometry}
\usepackage{parskip}
\setlength{\parskip}{6pt}
\setlength{\parindent}{0pt}

\usepackage{titlesec}
\usepackage{tocloft}
\usepackage{enumitem}

\usepackage[dvipsnames]{xcolor}
\usepackage{hyperref}
\usepackage{url}
\usepackage{bookmark}

\usepackage{fancyhdr}

\hypersetup{
  colorlinks=true,
  linkcolor=MidnightBlue,
  urlcolor=RoyalBlue,
  citecolor=ForestGreen,
  pdftitle={Pregled gradiva kursa},
  pdfpagemode=UseOutlines
}

\renewcommand{\contentsname}{Sadržaj}
\renewcommand{\cftsecleader}{\cftdotfill{\cftdotsep}}
\setlength{\cftbeforesecskip}{4pt}
\setlength{\cftbeforesubsecskip}{2pt}

\titleformat{\section}
  {\Large\bfseries\sffamily}
  {\thesection}{0.8em}{}
  [\vspace{0.3em}\titlerule]
\titleformat{\subsection}
  {\normalsize\bfseries}
  {\thesubsection}{0.6em}{}

\setlist[itemize]{itemsep=0.25em, topsep=0.3em, leftmargin=1.2cm}

\setlength{\headheight}{14.5pt}
\pagestyle{fancy}
\fancyhf{}
\fancyhead[L]{\small Računarstvo i društvo}
\fancyfoot[C]{\sffamily \thepage}
\renewcommand{\headrulewidth}{0.2pt}
\renewcommand{\footrulewidth}{0pt}

\raggedbottom
\clubpenalty=10000
\widowpenalty=10000

\title{%
\sffamily
\Large Univerzitet u Beogradu\\
\large Matematički fakultet\\[0.8em]
\textbf{Računarstvo i društvo}\\[1.2em]
\rule{\linewidth}{0.5pt}\\[0.6em]
\LARGE \textbf{Pregled gradiva}\\
\rule{\linewidth}{0.5pt}\\[1.2em]
}

\author{%
\sffamily
\begin{tabular}{@{}ll@{}}
\textbf{Student:} & Danilo Nikolaš (112/2021) \\
\textbf{Profesor:} & Dr. Sana Stojanović Durdević \\
\end{tabular}
}

\date{\sffamily Novembar 2025}

\begin{document}
\maketitle
\newpage
\tableofcontents
\newpage

\section{Istorijski pregled razvoja računarstva}
\subsection{Pregled gradiva}
Razvoj računarstva može se posmatrati kao niz postepenih promena, a ne kao jedna nagla revolucija. Njegovi počeci vezuju se za jednostavne mehaničke alatke kao što je abakus, a kasnije i prve kalkulatore koji su mogli obavljati osnovne operacije. Prva ideja o mašinama koje se mogu programirati pripisuje se Čarlsu Bebidžu, dok je Ada Lavlejs napravila korak dalje i opisala kako bi takva mašina mogla da izvršava unapred definisane instrukcije, što se danas smatra začetkom programiranja.

Prava prekretnica nastupa tokom 20. veka, sa pojavom tranzistora, a zatim i integrisanih kola i mikroprocesora. Zahvaljujući tome, računari su postali znatno manji, pouzdaniji i pristupačniji. U isto vreme, menjao se i način na koji se pišu programi, od bušenih kartica i kompleksnog asemblera, do jezika koji su postepeno postajali apstraktniji, modularni i prilagođeni ljudskom načinu razmišljanja. Ovo je posebno bilo izraženo sa pojavom objektno-orijentisanog pristupa.

Veliku ulogu u širenju tehnologije imao je i razvoj mreža. ARPANET je bio prvi korak ka umrežavanju računara na daljinu, dok su internet i kasnije Veb učinili informacije dostupnim svima i promenili način na koji ih koristimo.

Danas se računarstvo razvija u pravcu mobilnih uređaja, cloud servisa i \emph{edge} sistema, a istovremeno se intenzivno radi na veštačkoj inteligenciji i kvantnim računarima. Ipak, razvoj hardvera i dalje se najviše kreće ka tome da uređaji budu manji, brži i manje energetski zahtevni, kako bi bili praktičniji za svakodnevnu upotrebu.


\subsection{Predložene teme}
\begin{itemize}
  \item Evolucija računara od industrijske revolucije do mobilne ere \\
  \textbf{Link:} \url{https://www.computerhistory.org/timeline/computers/}
  \item Računarstvo kao pokretač globalizacije: kako je digitalna tehnologija promenila ekonomiju i društvo \\
  \textbf{Link:} \url{https://bastakiss.com/blog/our-blog-1/the-history-of-programming-language-standardization-and-why-it-matters-703}
\end{itemize}

\section{Etika}
\subsection{Pregled gradiva}
Etika u računarstvu bavi se pitanjem šta je ispravno, a šta problematično u svetu tehnologije. Nastala je iz klasične etike, ali su računari i internet otvorili nova pitanja, poput privatnosti, korišćenja podataka i odgovornosti onih koji prave softver. Razlog za razmišljanje o ovim pitanjima i dalje se nalazi u filozofskim pravcima poput deontologije, utilitarizma i etike vrline, koji na različite načine pokušavaju da objasne šta je moralno i do čega naše odluke dovode.

Razvojem računarskih sistema pojavila se potreba da se etička pitanja ne posmatraju samo teorijski, već i kroz praktične situacije, npr. kada je opravdano koristiti tuđe podatke, gde je granica digitalnog nadzora ili da li je etički učiti studente pisanju virusa u svrhu odbrane od njih. Pored toga, računarska etika se bavi i profesionalnim ponašanjem, kao i pravilima koja smanjuju rizik od zloupotrebe tehnologije.

Danas, etika u IT-u podrazumeva odgovornost prema korisnicima i društvu, promišljanje posledica koje tehnologija može izazvati i stalno preispitivanje načina na koji se sistemi razvijaju i koriste.

\subsection{Predložene teme}
\begin{itemize}
  \item Odgovornost algoritama i borba protiv pristrasnosti \\
  \textbf{Link:} \url{https://www.techtarget.com/searchsecurity/tip/How-AI-is-reshaping-threat-intelligence?int=off&pre=off&Offer=sy_lp10032025GOOGOTHR_GsidsSecurity_Akamai_KTO_IO335417_LI2952645&gad_campaignid=23075048613/}
  \item Profesionalna etika programera pod pritiskom rokova \\
  \textbf{Link:} \url{https://medium.com/@JohanneA/effective-deadlines-finishing-your-programming-projects-a366bb6e320e/}
\end{itemize}

\section{Upotreba interneta, spam i phishing}
\subsection{Pregled gradiva}
Internet je danas glavni način komunikacije i razmene informacija, ali je sa njegovom popularnošću porastao i broj različitih zloupotreba. Jedan od najčešćih problema je \emph{spam}, što predstavlja masovno slanje neželjenih poruka, najčešće u reklamne svrhe, ali često i kao prvi korak u ozbiljnije prevare. Posebno je rasprostranjen \emph{phishing}, gde se napadač lažno predstavlja, obično kao institucija ili servis od poverenja, kako bi naveo korisnika da otkrije lozinku, broj kartice ili druge privatne podatke.

Ove poruke najčešće deluju hitno ili zvanično kako bi izazvale brzu reakciju bez previše razmišljanja. Osim putem mejla, prevare se danas šire i kroz poruke na telefonima, društvenim mrežama, pa čak i linkovima koji vode do lažnih sajtova dizajniranih da izgledaju potpuno stvarno.

Zaštita od ovakvih prevara ne zavisi samo od tehnologije, već i od samih korisnika. Dobri filteri i dodatne sigurnosne mere pomažu, ali je ključno da ljudi znaju da prepoznaju sumnjive poruke pre nego što reaguju na njih.


\subsection{Predložene teme}
\begin{itemize}
  \item Psihologija phishing prevare i tehnike zaštite \\
  \textbf{Link:} \url{https://www.coursera.org/gb/articles/types-of-cyber-attacks?campaignid=20882109092&adgroupid=&device=c&keyword=&matchtype=&network=x&devicemodel=&creativeid=&assetgroupid=6485735763&targetid=&extensionid=&placement=&gad_campaignid=20875558740/}
  \item Evolucija antispam sistema i primena mašinskog učenja \\
  \textbf{Link:} \url{https://www.mimecast.com/content/anti-spam-software/}
\end{itemize}

\section{Cenzura i sloboda govora, deca i neprikladan sadržaj}
\subsection{Pregled gradiva}
Cenzura predstavlja kontrolu i ograničavanje informacija koje se smatraju neprikladnim, štetnim ili opasnim, a kroz istoriju je najčešće sprovode institucije sa velikim uticajem, kao što su država ili religijske organizacije. Ona može biti direktna, kada se sadržaj blokira ili zabranjuje, ili u obliku samocenzure, kada se pojedinci sami povlače i ustručavaju da iskažu mišljenje, najčešće iz straha od osude ili posledica.

Internet je dodatno promenio način na koji cenzura funkcioniše, informacije se šire brzo i teško ih je kontrolisati, ali i dalje se radi na tome da se ograniči pristup, bilo kroz zakone, blokade ili filtriranje sadržaja. U isto vreme, sloboda govora ostaje jedna od najvažnijih društvenih vrednosti, jer omogućava razmenu ideja, kritičko mišljenje i razvoj društva, ali postaje problem kada neko pokuša da odlučuje šta sme, a šta ne sme da se kaže.
Posebno je važno obratiti pažnju na decu, koja danas odrastaju uz internet i digitalne tehnologije. Internet im pruža znanje i mogućnosti, ali ih istovremeno izlaže sadržajima koji nisu za njihov uzrast. Zbog toga se koriste različite metode zaštite, poput softvera za filtriranje i nadzora sadržaja, ali i edukacije, jer tehnologija može da pomogne, ali ne i da u potpunosti zameni odrasle.

\subsection{Predložene teme}
\begin{itemize}
  \item Moderacija sadržaja i sloboda govora \\
  \textbf{Link:} \url{https://ecpr.eu/Events/Event/PanelDetails/8462/}
  \item Deepfake tehnologija i zaštita identiteta \\
  \textbf{Link:} \url{https://www.fortinet.com/resources/cyberglossary/deepfake/}
\end{itemize}

\section{Internet prevare i zavisnost od interneta}
\subsection{Pregled gradiva}
Internet nam je olakšao svakodnevni život, ali je istovremeno doneo i brojne prevare, od lažnih mejlova i profila na mrežama, do ozbiljnijih pokušaja zloupotrebe poverenja korisnika. Napadači često koriste poruke tako da da deluju hitno ili zvanično, uz sitne "nevidljive" promene (npr. jedno izmenjeno slovo u linku), koje je lako prevideti ako se ne obrati pažnja. Zato je važno biti pažljiv i proveravati izvore, označavati pouzdane adrese i dvaput razmisliti pre nego što se unesu osetljivi podaci.

Sa druge strane, internet je toliko prisutan u svakodnevnom životu da kod mnogih stvara potrebu da stalno budu online, što vremenom može da preraste u zavisnost i ostavi razne posledice. Kod mladih je to posebno vidljivo, jer društvene mreže i online igre često potisnu druženje uživo i mogu da utiču na fokus, san i odnose sa vršnjacima. Zato je važno na vreme prepoznati znakove zavisnosti i reagovati, bilo kroz razgovor, postavljanje granica u korišćenju interneta ili stručnu pomoć kada je potrebno.


\subsection{Predložene teme}
\begin{itemize}
  \item Uticaj micro-transakcija i loot-box mehanika u online igrama na razvoj zavisničkog ponašanja \\
  \textbf{Link:} \url{https://www.tuw.edu/psychology/psychology-behind-microtransactions/}
  \item Lažne recenzije i manipulacija potrošača \\
  \textbf{Link:} \url{https://pirg.org/edfund/resources/how-to-recognize-fake-online-reviews-of-products-and-services/}
\end{itemize}

\section{Bezbednost interneta}
\subsection{Pregled gradiva}
Sajber bezbednost obuhvata zaštitu računarskih sistema, podataka i korisnika od napada i zloupotreba. Posebno je važna zaštita lozinki: umesto čuvanja u običnom tekstu, koriste se heš funkcije i dodatni nasumični podatak "so" (eng. \emph{salt}), koji se dodaje lozinki pre heširanja, čime se otežava njeno otkrivanje. Preporučuje se i korišćenje višefaktorske autentifikacije, kao i upotrebe menadžera lozinki.

Bezbednost sistema unapređuje se tehnikama kao što su redovna ažuriranja, segmentacija mreže i primena pravila najmanjih privilegija. Ovi pristupi smanjuju \textbf{površinu napada}, tj. ukupan broj potencijalnih tačaka kroz koje napadač može dobiti pristup sistemu.

Najčešće savremene pretnje uključuju \emph{phishing} (prevare sa ciljem krađe podataka), \emph{ransomware} (zaključavanje podataka uz zahtev za otkup) i \emph{DDoS} napade (preopterećivanje servera). Zato dobra zaštita podrazumeva kombinaciju tehnologije, jasnih pravila bezbednosti i korisnika koji znaju kako da se zaštite od napada.

Praktičnu obuku pružaju CTF takmičenja (\emph{Capture The Flag}), gde učesnici rešavaju kontrolisane bezbednosne izazove i uče etičko hakovanje. U oblasti privatnosti, \emph{Privacy Sandbox} predstavlja alternativu kolačićima uz smanjenje praćenja korisnika, dok \emph{fingerprinting} označava identifikaciju uređaja na osnovu njegovih karakteristika. Razumevanje ovih tehnika važno je za očuvanje bezbednosti i privatnosti na internetu.


\subsection{Predložene teme}
\begin{itemize}
  \item Sigurnost lozinki i uloga menadžera lozinki \\
  \textbf{Link:} \url{https://www.waldenu.edu/programs/information-technology/resource/cybersecurity-101-why-choosing-a-secure-password-in-so-important/}
  \item Moderne metode autentifikacije (MFA, FIDO2, biometrija) \\
  \textbf{Link:} \url{https://www.loginradius.com/blog/identity/top-authentication-methods/}
\end{itemize}

\section{Privatnost na internetu i kolačići}
\subsection{Pregled gradiva}
Privatnost na internetu odnosi se na pravo korisnika da kontroliše ko ima pristup njegovim ličnim podacima i na koji način se oni koriste. Iako internet olakšava komunikaciju i pristup informacijama, ostavlja i digitalne tragove, ponekad zato što to sami dozvolimo, a ponekad bez našeg znanja.

Veliki deo podataka danas se prikuplja kroz kolačiće, male fajlove koji pamte jezik, podešavanja, korpu za kupovinu, ali i prate ponašanje korisnika na sajtu. Oni mogu biti kratkotrajni ili trajni, a posebno su diskutabilni kolačići treće strane, koji omogućavaju praćenje aktivnosti korisnika i van sajta na kom su nastali. Zbog zloupotreba i netransparentnog prikupljanja podataka, danas je obavezno obavestiti korisnika i dobiti njegov pristanak pre njihovog korišćenja

Iako zakoni pokušavaju da zaštite privatnost, korisnici često ne čitaju politike korišćenja i nisu u potpunosti svesni šta prihvataju. Zato je važno koristiti zaštitne mehanizme poput privatnog režima pretraživača, brisanja i blokiranja kolačića, enkripcije i VPN servisa. Možda i najbitnije je da korisnik shvati rizike deljenja svojih podataka i nauči kako da ostavlje svoje podatke online.

\subsection{Predložene teme}
\begin{itemize}
  \item Budućnost interneta bez kolačića trećih strana \\
  \textbf{Link:} \url{https://ignitevisibility.com/web-without-third-party-cookies/}
  \item Manipulativni dizajn (dark patterns) u cookie banerima \\
  \textbf{Link:} \url{https://www.nsw.gov.au/departments-and-agencies/fair-trading/dark-patterns/}
\end{itemize}

\end{document}
